% =========================================================================
% CAVE-HAR — 3DV 2027 draft (target deadline: Aug 28, 2026 AoE)
% -------------------------------------------------------------------------
% NOTE: 3DV uses the CVF-style two-column format. Before submission,
% replace this preamble with the official 3DV 2027 author kit (expected on
% https://3dvconf.github.io/2027/call-for-papers/) and delete the \todo
% macro output by switching \todofalse. All open gaps are marked \todo{...}
% and indexed in ../GAPS.md with owner + producing phase.
% =========================================================================
\documentclass[10pt,twocolumn,letterpaper]{article}

\usepackage[margin=1in]{geometry}   % replaced by official style file
\usepackage{times}
\usepackage{graphicx}
\usepackage{booktabs}
\usepackage{multirow}
\usepackage{amsmath}
\usepackage{pifont}   % \ding crosses in the dataset comparison table
\usepackage{xcolor}
\usepackage[pagebackref,breaklinks,colorlinks,citecolor=blue]{hyperref}

% ---- draft machinery ----------------------------------------------------
\newif\iftodos \todostrue   % \todosfalse hides all TODO markers
\newcommand{\todo}[1]{\iftodos\textcolor{red}{\textbf{[TODO: #1]}}\fi}
\newcommand{\DATASET}{CAVE-HAR\todo{数据集正式命名待定 (G1)}}
\newcommand{\numclips}{\todo{清洗后片段总数 (G8)}}

\begin{document}

\title{\DATASET: A Multi-View Benchmark for Human Action Recognition\\
in Immersive CAVE Environments}

\author{Anonymous 3DV 2027 submission\\Paper ID \todo{****}}
\date{}
\maketitle

% =========================================================================
\begin{abstract}
Human action recognition (HAR) inside large-scale immersive systems such as
CAVEs (Cave Automatic Virtual Environments) promises marker-less
interaction, automatic procedure assessment, and behavioral analysis for
safety-critical training. Yet every widely used HAR benchmark implicitly
assumes conditions that a CAVE violates: frontal tripod-height viewpoints,
stable photometric backgrounds, and offline evaluation. We present
\DATASET{}, the first HAR benchmark captured inside an operational
projection-based CAVE. It provides synchronized multi-view RGB video from
seven cameras---including the top-mounted, steeply pitched views forced by
CAVE geometry---together with optical-motion-capture-grade 3D skeleton
ground truth (SMPL layout, 22 joints) for 15 actors performing 7 general
and 11 VR-interaction actions across 6 projected scenes. Because the
projected VR content changes the entire visual background under controlled
conditions, \DATASET{} additionally enables a background-shift evaluation
that is impossible to stage in real-world scenes. We define four evaluation
protocols (cross-subject, cross-view, cross-scene, and cross-domain from
NTU~RGB+D~120), plus a latency-aware protocol that reports computation and
observation latency for streaming recognition. Benchmarking
\todo{N (G13)} representative skeleton, video, and multimodal models
reveals \todo{核心发现一句话, 需实验后填 (G13)}. The dataset, evaluation
code, and correction toolchain will be publicly released.
\end{abstract}

% =========================================================================
\section{Introduction}
Projection-based immersive systems such as the CAVE~\cite{cruzneira1993cave}
remain the platform of choice when virtual reality must be combined with
free full-body movement, unobstructed vision of one's own body, and
multi-user co-presence---requirements typical of industrial safety
training, rehabilitation, and simulation. Human action recognition inside
such environments would enable automatic evaluation of trainee procedures,
marker-less interaction, and quantitative behavioral analysis.

Modern HAR models are, however, trained and evaluated almost exclusively on
datasets that encode three assumptions. First, \emph{fixed frontal
geometry}: cameras in NTU~RGB+D~\cite{shahroudy2016ntu,liu2020ntu120} and
its successors observe subjects from tripod height at moderate distances,
whereas a CAVE forces cameras to its ceiling rim, producing steep pitch
angles and systematic self-occlusion of hands and head during interaction.
Second, \emph{appearance stability}: laboratory backgrounds are static,
whereas a CAVE's walls and floor \emph{are} the display---the background
changes radically with the projected scene and its light spills onto the
actor. Third, \emph{offline evaluation}: accuracy is reported on trimmed
clips without latency constraints, whereas interactive VR imposes tight
motion-to-response budgets. The gap between these assumptions and CAVE
reality is unquantified because no suitable benchmark exists.

We close this gap with \DATASET{}, captured in a
\todo{CAVE 尺寸/面数确认: 幻灯片同时出现 ``5 faces active'' 与 ``6-sided''
(G2)} projection CAVE instrumented with 12 OptiTrack Prime$^{\mathrm{x}}$13
motion-capture cameras and 7 RGB cameras. Our contributions are:
\begin{itemize}
  \item \textbf{Dataset.} \numclips{} clips of 15 actors (9M/6F), 6
        projected scenes (3 indoor, 3 outdoor), 7 general + 11
        VR-interaction actions, with synchronized 7-view RGB and
        mocap-grade SMPL-22 skeletons---to our knowledge the first HAR
        corpus recorded inside an operational CAVE, and the first with
        top-mounted viewpoints as a primary condition.
  \item \textbf{Human-verified ground truth.} A correction pipeline
        combining per-clip quality scoring, two-view triangulated editing,
        and bone-length-preserving inverse-kinematics tools; all released
        test-split skeletons are manually verified
        (\todo{修正帧占比 X\% (G9)}, inter-annotator agreement
        \todo{Y mm (G10)}).
  \item \textbf{Protocols.} Cross-subject, cross-view, cross-scene, and
        NTU$\rightarrow$\DATASET{} cross-domain splits, plus a
        latency-aware streaming protocol reporting computation and
        observation latency.
  \item \textbf{Baselines.} \todo{N (G13)} models spanning skeleton GCNs
        and transformers, video networks, and multimodal fusion, with an
        analysis of where and why they fail under CAVE conditions.
\end{itemize}

% =========================================================================
\section{Related Work}
\paragraph{HAR datasets.}
NTU~RGB+D~\cite{shahroudy2016ntu} and NTU~RGB+D~120~\cite{liu2020ntu120}
established the dominant evaluation setting: 106 subjects, 120 classes,
Kinect-estimated skeletons, frontal/side views. ANUBIS~\cite{yu2022anubis}
added rear viewpoints and finer-grained interactions with Azure-Kinect
skeletons. Kinetics~\cite{kay2017kinetics} scales to web video without 3D
ground truth. All estimate skeletons from consumer depth sensors or 2D
pose; none offers optical-mocap-grade joints, top-mounted views, or
projected-background variation. Our comparison appears in
Table~\ref{tab:datasets}.

\paragraph{User behavior corpora in VR.}
LocoVR~\cite{takeyama2024locovr} records multi-user locomotion trajectories
in HMD-based VR; HUMOTO~\cite{lu2025humoto} provides mocap human--object
interactions without multi-view RGB. HMD-based corpora observe users
wearing headsets that occlude the face and alter motion; CAVE users wear
only lightweight shutter glasses, keeping appearance close to natural while
introducing the projection-specific challenges we target. RoCoG-v2
\cite{reddy2023rocog} studies synthetic-to-real domain shift for gestures,
methodologically related to our cross-domain protocol.

\paragraph{Models.}
Skeleton-based recognition evolved from graph convolutions
(ST-GCN~\cite{yan2018stgcn}, CTR-GCN~\cite{chen2021ctrgcn}) to
attention and prototype-based designs (SkateFormer~\cite{do2024skateformer},
ProtoGCN~\cite{liu2025protogcn}), state-space models
(SkelMamba~\cite{martinel2024skelmamba}), and pretrained motion encoders
(MotionBERT~\cite{zhu2023motionbert}). Video models progressed from
two-pathway CNNs (SlowFast~\cite{feichtenhofer2019slowfast}) to
transformers (Video Swin~\cite{liu2022videoswin},
VideoMAE~v2~\cite{wang2023videomae2}) and state-space architectures
(VideoMamba~\cite{li2024videomamba}); language--vision pretraining enables
zero-shot transfer (ActionCLIP~\cite{wang2021actionclip}); and
PoseConv3D~\cite{duan2022poseconv3d} fuses pose heatmaps with RGB. We
benchmark representatives of each family.

\paragraph{Online and latency-aware evaluation.}
Online action detection~\cite{degeest2016oad} and streaming perception
argue that offline accuracy overstates deployable performance; recent work
separates \emph{computation} latency from \emph{observation} latency
\cite{eyzaguirre2024squid}. VR interaction makes these constraints
first-class, motivating our protocol.

% =========================================================================
\section{The \DATASET{} Dataset}

\subsection{Capture environment}
Recording took place in the \todo{机构匿名化处理后的 CAVE 描述 + 尺寸
(G2)} CAVE. Twelve OptiTrack Prime$^{\mathrm{x}}$13 cameras capture 3D
motion at 240\,Hz; seven RGB cameras (\todo{型号/分辨率/帧率 (G3)},
nominally 30\,fps) cover the volume in a $3{\times}2$ grid on the left,
center, and right rim at high and low mounting positions, plus one
diagonal rear view. Fig.~\ref{fig:rig}\todo{相机布置图 (G4)} shows the
layout; the ``high'' row is top-mounted at pitch angles of
\todo{俯仰角度数 (G3)}, the geometry that distinguishes CAVE deployments
from tripod-based datasets.

Two collection rounds were performed. Round~1 (portal-configuration LED
CAVE, \todo{7 或 8 (幻灯片不一致, G5)} mocap + 2 RGB cameras) served as a
pilot; Round~2, described above, is the release configuration. Actors wore
everyday clothing with 37 retro-reflective markers; Round~2 used low-profile
marker stickers to minimize visual interference with the RGB streams.

\subsection{Actors, scenes, and actions}
Fifteen actors (9 male, 6 female; \todo{年龄范围/身高范围, 需 roster (G6)})
each performed in six projected scenes (three indoor, three outdoor),
approximately 20 effective minutes per scene and roughly 41{,}400 mocap
frames per actor. The action vocabulary (Table~\ref{tab:actions}) contains
seven general actions shared with existing benchmarks---enabling direct
cross-domain comparison---and eleven VR-interaction actions elicited by the
projected applications, from controller locomotion to catching virtual
fish with a physical net prop.

\begin{table}[t]\centering\small
\caption{Action vocabulary: 7 general + 11 VR-interaction classes.}
\label{tab:actions}
\begin{tabular}{@{}ll@{}}
\toprule
\textbf{General} & \textbf{VR-interaction} \\
\midrule
walking      & move via controller \\
running      & catch fish with net \\
jumping      & grab box \\
bending down & measure length \\
standing     & wave \\
squatting    & throw \\
raising hand & cut \\
             & shoot \\
             & pick up item \\
             & go up stairs \\
             & dodge car \\
\bottomrule
\end{tabular}
\end{table}

\begin{table*}[t]\centering\small
\caption{Comparison with existing HAR and VR-behavior datasets.
\todo{逐格核对数字并补齐空格 (G7)}}
\label{tab:datasets}
\begin{tabular}{@{}lcccccc@{}}
\toprule
Dataset & Subjects & Classes & Clips & Skeleton source & Top-mounted views & Controlled background shift \\
\midrule
NTU RGB+D 120~\cite{liu2020ntu120} & 106 & 120 & 114{,}480 & Kinect v2 & \ding{55} & \ding{55} \\
ANUBIS~\cite{yu2022anubis} & 80 & 102 & 66{,}232 & Azure Kinect & \ding{55} & \ding{55} \\
LocoVR~\cite{takeyama2024locovr} & \todo{} & locomotion & \todo{} & HMD tracking & \ding{55} & \ding{55} \\
\DATASET{} (ours) & 15 & 18 & \numclips & optical mocap + MoSh++ & \checkmark & \checkmark \\
\bottomrule
\end{tabular}
\end{table*}

\subsection{Calibration and synchronization}
RGB intrinsics and extrinsics follow Zhang's method~\cite{zhang1999calib}
with per-session refinement; extrinsic drift is monitored by
re-projecting triangulated 2D poses and gating on median residual
(\todo{标定精度: 平均重投影误差 px (G11)}). Lacking hardware genlock,
video--skeleton alignment uses a per-clip timeline established from the
mocap clock, refined per camera with sub-frame offsets and verified by
overlay inspection; final alignment accuracy is \todo{同步精度 ±帧 (G11)}.

\subsection{Ground-truth construction}
Raw 240\,Hz markers are solved to SMPL~\cite{loper2015smpl} bodies with
MoSh++~\cite{mahmood2019amass}; we release the 22-joint SMPL subset
(hand joints excluded owing to sparse hand markers) in both world
coordinates and pelvis-centered normalized form. Sparse marker coverage
and interaction props cause characteristic failures---frozen or drifting
limb segments---which we repair with a purpose-built toolchain: (i) every
clip receives an automatic quality score combining bone-length deviation,
robust velocity-spike detection, frozen-joint-while-moving runs, and
missing-marker fraction, which prioritizes review worst-first;
(ii) annotators correct joints by dragging in two calibrated views
simultaneously (exact triangulated depth) with keyframe--spline
propagation; and (iii) a bone-length-preserving two-bone IK mode solves
whole-limb placements from a single wrist/ankle drag, guaranteeing
physically plausible corrections that never fold a straight limb. All
test-split clips are fully verified; training-split clips are corrected
worst-first by quality score with per-frame edit masks released
(\todo{银标准修正比例 (G9)}). Corrections modified \todo{X\% (G9)} of
test-split frames; dual-annotation of a 5\% sample shows mean joint
disagreement of \todo{Y mm (G10)}.

\subsection{Statistics}
\todo{最终统计表: 每类样本数/时长分布/marker 丢失率/QC 分布直方图,
清洗后生成 (G8)}

% =========================================================================
\section{Benchmark Protocols}
\paragraph{Cross-subject (X-Sub, primary).}
Three fixed 10-train/5-test actor partitions (seeded, gender-balanced
\todo{性别核对后确认 (G12)}); we report mean$\pm$std over the three
partitions to control for small-cohort variance.
\paragraph{Cross-view (X-View).}
Train on low-rim views, test on top-mounted and diagonal-rear views,
isolating the viewpoint gap.
\paragraph{Cross-scene (X-Scene).}
Train on the three indoor projections, test on the three outdoor ones:
identical actions, radically different projected backgrounds---a
controlled background-shift experiment unavailable in real-world corpora.
\paragraph{Cross-domain.}
Models pretrained on NTU RGB+D 120 are evaluated zero-shot and after
fine-tuning, quantifying the lab-to-CAVE domain gap that motivates this
work.
\paragraph{Metrics.}
Top-1/Top-5 accuracy and mean per-class accuracy (the 18-class vocabulary
is imbalanced across interaction types); confusion analyses separate
general from interaction classes.
\paragraph{Latency-aware protocol.}
For streaming deployment we report (i)~\emph{computation latency}:
single-clip inference time on one \todo{评测 GPU 型号 (G14)}, and
(ii)~\emph{observation latency}: the minimum observed-frame budget at
which a model reaches 90\% of its full-clip accuracy
\cite{degeest2016oad,eyzaguirre2024squid}, yielding accuracy--latency
curves (Fig.~\todo{latency 图 (G15)}).

% =========================================================================
\section{Experiments}
\todo{整节数字待 P3 实验产出 (G13). 表格骨架已按分工排好.}

\begin{table}[t]\centering\small
\caption{Baselines on \DATASET{} (Top-1 / mean per-class, \%).
\todo{全部数字 (G13)}}
\label{tab:baselines}
\begin{tabular}{@{}llccc@{}}
\toprule
Model & Modality & X-Sub & X-View & X-Scene \\
\midrule
ST-GCN~\cite{yan2018stgcn}             & Skel. & \todo{} & \todo{} & \todo{} \\
CTR-GCN~\cite{chen2021ctrgcn}          & Skel. & \todo{} & \todo{} & \todo{} \\
SkateFormer~\cite{do2024skateformer}   & Skel. & \todo{} & \todo{} & \todo{} \\
SkelMamba~\cite{martinel2024skelmamba} & Skel. & \todo{} & \todo{} & \todo{} \\
MotionBERT~\cite{zhu2023motionbert}    & Skel. & \todo{} & \todo{} & \todo{} \\
SlowFast~\cite{feichtenhofer2019slowfast} & RGB & \todo{} & \todo{} & \todo{} \\
Video Swin~\cite{liu2022videoswin}     & RGB   & \todo{} & \todo{} & \todo{} \\
VideoMAE v2~\cite{wang2023videomae2}   & RGB   & \todo{} & \todo{} & \todo{} \\
VideoMamba~\cite{li2024videomamba}     & RGB   & \todo{} & \todo{} & \todo{} \\
ActionCLIP~\cite{wang2021actionclip}   & RGB+T & \todo{} & \todo{} & \todo{} \\
PoseConv3D~\cite{duan2022poseconv3d}   & Skel.+RGB & \todo{} & \todo{} & \todo{} \\
\bottomrule
\end{tabular}
\end{table}

\paragraph{Cross-domain gap.} \todo{NTU 预训练 zero-shot/finetune 表 (G13)}
\paragraph{Viewpoint analysis.} \todo{逐相机准确率 + 遮挡消融 (G13)}
\paragraph{Zero-shot recognition.} \todo{ActionCLIP 零样本结果 (G13)}
\paragraph{Latency.} \todo{accuracy-latency 曲线与讨论 (G13, G15)}
\paragraph{Qualitative results.} \todo{成功/失败案例可视化 (G16)}

% =========================================================================
\section{Limitations and Ethics}
\DATASET{} is captured in a single CAVE installation with 15 actors; the
cohort size follows from mocap-suit logistics and is mitigated by the
multi-partition X-Sub protocol, but broader demographic coverage is future
work. Hand joints are excluded from the official set because sparse hand
markers make them unreliable. Training-split skeletons are corrected
worst-first rather than exhaustively; per-frame edit masks make the
verification status of every joint transparent. All participants gave
written informed consent covering public release of RGB video and derived
data under an institution-approved protocol
(\todo{伦理批号, 投稿系统需要 (G17)}); raw uncorrected captures are not
distributed. \todo{面部是否模糊处理的最终决定 (G17)}

\section{Conclusion}
\DATASET{} turns an operational CAVE into a benchmark, pairing
mocap-grade skeletons with the viewpoints, backgrounds, and latency
constraints that immersive deployment actually imposes. Our baselines show
\todo{一句话结论 (G13)}, and the released toolchain lowers the cost of
extending the corpus. We hope \DATASET{} seeds a line of HAR research that
treats immersive environments as a first-class setting rather than an
afterthought.

{\small
\bibliographystyle{ieee_fullname}
\bibliography{refs}
}

\end{document}
